\documentclass[10pt,a4paper]{amsart}
\usepackage[margin=19mm]{geometry}
\usepackage{amsmath,amssymb,mathtools}
\usepackage[scr=rsfs,cal=euler]{mathalfa}
\usepackage{xfrac}
\usepackage[unicode,hidelinks,bookmarks=false]{hyperref}
\newcommand{\ie}{i.e.\ }
\newcommand{\cf}{cf.\ }
\newcommand{\resp}{respectively\ }
\newcommand{\Subsection}{\subsection}
\providecommand{\nobreakdash}{\nobreak\hbox{-}}
\setlength{\emergencystretch}{2em}
\multlinegap0pt
\usepackage{amssymb}
\newtheorem{theorem}{Theorem}
\title{Stability of Partitions Induced by Nearest-Center Assignment Under Perturbations}
\author{MD Nahidul Hasan Sabit, Faija Anjum}
\date{Source: arXiv:2604.15578v1 (CC BY 4.0)}
\begin{document}
\maketitle

\noindent\emph{Source and licence.} Stability of Partitions Induced by Nearest-Center Assignment Under Perturbations. Version arXiv:2604.15578v1 (CC BY 4.0), distributed by arXiv under the Creative Commons Attribution 4.0 International licence, http://creativecommons.org/licenses/by/4.0/, as recorded in the arXiv licence metadata for this exact version and shown on the abstract page https://arxiv.org/abs/2604.15578v1. Full licence notice: \texttt{licenses/arxiv-2604.15578-CC-BY-4.0.txt}.\par\smallskip

\noindent\textit{Editorial scope.} Counted unit: the article's theorem theorem (source0.tex lines 157--170). The article's complete original proof of it is delivered with it (source0.tex lines 172--208, one \texttt{proof} environment of 37 lines). No mathematics is written, repaired or re-derived: the statement and the proof are the article's own lines, in the article's own notation, and no retained line is substituted or re-flowed.  No further source-local context is needed: every cross-reference of every retained line resolves inside the two delivered spans.  Nothing was omitted from any retained span, so the package carries no omission record.  No retained line contains a citation, so the wrapper carries no bibliography and no bibliography entry.  No retained line contains a figure environment, an image-inclusion command or drawing code, so no image is delivered and the package declares no asset dependency.  Editorial formatting only, in addition: an AMS \texttt{amsart} preamble that restates the article's own declarations and macros used by the retained lines, namely the environments \texttt{theorem} on separate counters, together with the standard packages those lines need.  The retained environments print in this document as Theorem 1 in order of appearance, whereas the article prints the same items with the numbering of its own full document.  The complete native source of the article as distributed in the arXiv e-print for this exact version is bundled unchanged as \texttt{sources/198621/source0.tex}.
\begin{theorem}
Let
\[
X' = (x_1', x_2', \ldots, x_n')
\]
be a perturbation of $X$ such that
\[
\|X - X'\| \le \varepsilon.
\]
If
$\varepsilon < \frac{\gamma_{\min}}{2},$
no point changes its assigned center under the perturbation. Consequently,
$A(X') = A(X).$
\end{theorem}

\begin{proof}
Fix $i \in \{1, \ldots, n\}$ and let $x_i'$ be the perturbed point. By definition of the perturbation size,
$\|x_i - x_i'\| \le \varepsilon.$
We compare the distances from $x_i'$ to the centers. By the triangle inequality, for $j = 1,2$,
\[
\|x_i' - c_j\| \le \|x_i - c_j\| + \varepsilon, \qquad
\|x_i' - c_j\| \ge \|x_i - c_j\| - \varepsilon.
\]

Assume that $x_i$ is assigned to $c_1$, so that
\[
\|x_i - c_1\| \le \|x_i - c_2\|.
\]
Then
\[
\|x_i - c_2\| - \|x_i - c_1\| = \gamma_i \ge \gamma_{\min}.
\]

For the perturbed point, we have
\[
\|x_i' - c_2\| - \|x_i' - c_1\|
\ge (\|x_i - c_2\| - \varepsilon) - (\|x_i - c_1\| + \varepsilon)
= \gamma_i - 2\varepsilon.
\]

Since $\varepsilon < \gamma_{\min}/2 \le \gamma_i/2$, we obtain
\[
\gamma_i - 2\varepsilon > 0.
\]
Thus,
\[
\|x_i' - c_2\| > \|x_i' - c_1\|,
\]
so $x_i'$ remains assigned to $c_1$.

The same argument applies if $x_i$ is assigned to $c_2$. Therefore, no point changes its assigned center under the perturbation. Since cluster membership is determined by assigned centers and no assignments change, all pairwise relations of belonging to the same cluster remain unchanged. Thus, the induced partition does not change, and $A(X') = A(X)$.
\end{proof}

\end{document}
